\section{Introduction}
\label{sec:introduction}

An action-chunk policy predicts commands for several future control steps from
each observation~\cite{zhao2023act}. Its executor must then decide how long or
how repeatedly to use those predictions before requesting a new chunk. Common
executors expose one temporal setting for the full action vector, although the
vector combines translation, rotation, and gripper commands. A single setting
is behaviorally neutral only if these components respond similarly when their
predictions become stale.

The evaluated ACT policy shows a striking component structure on the Object
development cohort. At a source age of 1.00~s, success was 8.73\% with stale
translation, 9.52\% with the full arm stale, 42.06\% with stale rotation,
44.44\% with all components fresh, and 64.29\% with a stale gripper. These
exactly block-matched values combine audited historical anchors with later
reviewer-directed component probes; they are a comparative characterization,
not one preregistered five-condition factorial. Staling translation alone
nearly reproduces the degradation seen when the full arm is stale. Rotation
staleness incurred little detectable cost relative to Fresh on this cohort
($-2.38$ points, $p=0.607$, with both intervals including zero). The magnitude
of the stale-gripper simple effect varies substantially across cohorts, so the
64.29\% value is not evidence of a general benefit.

We measure this structure with a same-target diagnostic. At control step $t$,
one component can use offset 0 from the chunk queried at $t$, while another can
use offset 20 from the chunk queried 20 steps earlier. Both entries were
predicted for physical step $t$. More generally, query time $q$ and chunk
offset $k$ obey $q+k=t$. Holding the target fixed therefore couples an older
source observation to a farther-ahead chunk prediction. We measure their joint
effect and do not separate source age from prediction lookahead. Every probe
queries the policy once per executed step, making it a diagnostic intervention
rather than a deployment executor.

This study makes three contributions. First, a preregistered same-target test
shows a large arm--gripper asymmetry, and a separate development sweep preserves
its ordering over source ages from 0.10 to 1.60~s. Second, component probes show
strong translation--rotation heterogeneity within the arm and reveal that
same-target prediction dispersion ranks these components in the wrong order.
Third, executable component-resolved schedules improve success when gripper
commitment is extended at two fixed arm replanning cadences, while coherent H16
remains the strongest overall operating point.

The claims are deliberately bounded. The same-target results diagnose temporal
source assignment under a policy-query rate of 1, whereas deployment schedules
replan less often. The central behavioral evidence concerns the evaluated ACT
policies in LIBERO~\cite{liu2023libero}; it does not establish cross-policy
generality. Coherent H16 remains strongest in the frozen static benchmark, and
the diagnostics tested here do not identify the mechanism behind the component
ordering.
